\documentclass[12pt,a4paper]{article}
\usepackage[utf8]{inputenc}
\usepackage{amsmath,amssymb}
\usepackage{graphicx}
\usepackage{hyperref}
\usepackage{geometry}
\geometry{margin=2.5cm}

\begin{document}

\section{Theoretical Background}

This chapter presents the theoretical foundation underpinning the design and implementation of the proposed system. It begins with a broad overview of artificial intelligence and its applications in education, and then progressively narrows its focus to the specific technologies employed in this project: Large Language Models, Retrieval-Augmented Generation, Multi-Agent AI architecture with Human-in-the-loop mechanisms, document processing and vector databases, prompt engineering, and web development technologies. For each technology, the general principles are first established, followed by an explanation of how it is applied within the context of the exam generation system.

% ------------------------------------------------------------
\subsection{Overview of Artificial Intelligence}

Artificial Intelligence (AI) is a field of study concerned with the development of computer systems capable of performing tasks that normally require human intelligence, such as learning, reasoning, pattern recognition, language understanding, and decision-making. The overarching goal of AI is to build systems that can perceive information, process it, and generate outputs appropriate to a given problem.

The applications of AI span numerous domains, including healthcare, finance, e-commerce, and education. In the educational sector, AI serves two primary purposes. First, it supports the automation of administrative and management tasks, thereby reducing the workload on educators. Second, and more importantly, it contributes directly to pedagogical activities through content personalization, adaptive learning pathways, automated assessment, and intelligent content generation. This dual role demonstrates the significant potential of AI in enhancing both the efficiency and the quality of educational practice.

Among the most impactful directions in modern AI is Generative AI, particularly Large Language Models (LLMs), which are capable of generating coherent text, summarizing content, answering questions, and constructing knowledge from input data. However, for reliable use in academic contexts, AI systems must be designed with controllability and safety in mind, ensuring that outputs are factually grounded, contextually appropriate, and subject to appropriate human oversight.

The technologies explored in the subsequent sections — LLM, RAG, Multi-Agent AI, Human-in-the-loop, and Prompt Engineering — collectively address these requirements. They form the technological stack on which the proposed exam generation system is built, enabling a process that is both automated and well-controlled.

% ------------------------------------------------------------
\subsection{Large Language Models}

Large Language Models (LLMs) are artificial intelligence models trained on massive corpora of textual data to learn the patterns, structures, and relationships present in natural language. Built upon the Transformer architecture introduced by Vaswani et al. (2017), LLMs use self-attention mechanisms to process input text by relating every token to every other token in the sequence. This allows the model to capture long-range dependencies and contextual relationships within text — a capability that was difficult to achieve with earlier recurrent architectures.

Given a sequence of text as input, an LLM predicts the next token by estimating a probability distribution over its vocabulary, conditioned on the entire preceding context. Through large-scale pre-training followed by fine-tuning, the model learns to represent linguistic meaning, model conceptual relationships, and produce contextually appropriate responses. This mechanism underlies the model's ability to perform a broad spectrum of tasks including text generation, question answering, summarization, semantic analysis, and structured content creation.

The development of LLMs has marked a major turning point in natural language processing, substantially expanding the practical applications of AI across domains. In educational technology, LLMs are particularly valuable for tasks such as exam blueprint construction, question generation across various formats and cognitive levels, and automated content quality assessment. However, because LLMs operate through probabilistic prediction based on knowledge acquired during training, they are prone to generating inaccurate, hallucinated, or contextually misaligned content. In high-stakes educational contexts, these limitations must be mitigated through grounded retrieval, structured generation pipelines, and human supervision — a strategy adopted throughout the design of the proposed system.

% ------------------------------------------------------------
\subsection{Retrieval-Augmented Generation}

Retrieval-Augmented Generation (RAG) is an architectural pattern that augments the generation capabilities of LLMs with information retrieved from external knowledge sources. Rather than relying solely on parametric knowledge encoded during training, a RAG system retrieves relevant document segments at inference time and incorporates them into the model's context window. This approach produces outputs that are more factually grounded, more closely aligned with the provided source materials, and significantly less likely to contain hallucinated content.

A typical RAG workflow consists of two primary stages. The first stage is retrieval: given an input query, the system converts both the query and the stored knowledge segments into dense vector representations using an embedding model. It then searches the vector database for the most semantically similar segments, typically using cosine similarity or dot-product scoring. A reranking step may optionally be applied to further prioritize the most contextually relevant results before they are passed to the generator. The second stage is generation: the LLM receives both the original query and the top retrieved segments as context, and generates a response grounded in the retrieved information.

In educational applications, RAG is especially valuable because it ensures that generated content is derived from the specific teaching materials selected by the instructor rather than from the model's general knowledge. This improves the topical relevance of generated questions, constrains outputs to the intended knowledge scope, and provides a verifiable basis for assessing the alignment between examination content and curriculum. The retrieved segments also serve as source evidence that can be displayed alongside generated questions, supporting transparency and quality review.

In this project, RAG is embedded at multiple stages of the exam generation pipeline. After teaching documents are processed and stored in a retrieval-ready form, the system queries relevant knowledge segments based on the scope selected by the teacher. These segments are then used as context for blueprint construction, question generation, and quality validation, ensuring that the resulting examination is both document-grounded and scope-consistent.

% ------------------------------------------------------------
\subsection{Multi-Agent AI and Human-in-the-Loop}

% 1.4.1
\subsubsection{Multi-Agent AI}

Multi-Agent AI is an architectural approach in which a complex task is decomposed into subtasks, each handled by a specialized intelligent agent. Rather than relying on a single monolithic model to perform the entire workflow, this organization assigns distinct responsibilities — such as knowledge retrieval, planning, content generation, and quality validation — to separate agents that communicate through a defined protocol. This decomposition improves modularity, observability, and the ability to detect and correct errors at the appropriate stage of the pipeline.

In LLM-based systems, agents are typically implemented as reasoning-and-generation loops guided by a system prompt and supported by external tools or data sources. Communication between agents may follow a sequential pattern, in which the output of one agent becomes the input of the next, or a more dynamic pattern using shared state and conditional routing. The use of a shared message protocol — such as that provided by the LangGraph framework employed in this project — allows agents to maintain a structured conversation history, enabling stateful coordination across pipeline stages.

For the task of automated exam generation, the multi-agent approach is particularly well-suited because the workflow inherently involves distinct phases: understanding the curriculum scope, constructing an exam blueprint, generating questions at appropriate cognitive levels, validating output quality, and incorporating human feedback. Assigning each phase to a dedicated agent enables focused optimization of individual components while preserving a clear overall structure.

In this project, the exam generation pipeline is organized into specialized agents — including the retrieval agent, outline agent, builder agent, clarification agent, and orchestrator — coordinated through a directed graph. This hierarchical and stateful organization enhances the system's controllability, supports granular tracking of generation progress, and enables targeted intervention at any stage of the process.

% 1.4.2
\subsubsection{Human-in-the-Loop}

Human-in-the-loop (HITL) is a design pattern that integrates human participation into critical stages of an AI pipeline. Rather than operating as a fully autonomous system, a HITL-enabled pipeline exposes one or more checkpoints at which a human reviews, evaluates, approves, or revises the system's intermediate or final output. This pattern is essential for applications where correctness, pedagogical appropriateness, and contextual suitability cannot be guaranteed by the model alone.

In the context of educational content generation, the stakes are particularly high: a poorly constructed examination can misrepresent student knowledge, reinforce misconceptions, or fail to assess the intended learning outcomes. HITL addresses this risk by ensuring that a qualified educator retains meaningful control over the final output. At each checkpoint, the teacher can verify factual accuracy, confirm alignment with the curriculum, adjust question framing or difficulty, and ultimately approve or request revision of the generated content.

In this project, HITL checkpoints are embedded at two key stages of the exam generation pipeline. At the first checkpoint, the teacher reviews and approves or rejects the generated exam blueprint — including scope coverage and Bloom's taxonomy distribution — before question generation begins. At the second checkpoint, the teacher reviews individual questions, edits content directly, submits AI-assisted revision requests, and makes the final approval decision. This two-stage structure combines the efficiency of AI-driven automation with the accountability and domain expertise of the instructor.

% ------------------------------------------------------------
\subsection{Document Processing and Vector Database}

% 1.5.1
\subsubsection{Document Processing}

Document processing is the preparatory stage that transforms raw teaching materials into a structured, retrieval-ready format. Teaching documents typically exist in heterogeneous formats — such as PDF, DOCX, or PPTX — and contain a mixture of text, headings, figures, tables, and formatting that is not directly amenable to semantic search or content generation. The document processing pipeline therefore performs a series of transformation steps to extract clean text, recover structural information, and produce semantically coherent knowledge units.

The first step is text extraction, which involves parsing the document format to recover the raw textual content. For PDF documents, this may rely on heuristic-based heading detection using visual cues such as font size, bold weight, and spatial positioning, as well as layout analysis to handle multi-column formats and page headers or footers. For DOCX and PPTX files, the native structural elements (headings, paragraphs, lists) can often be extracted directly through the document object's internal XML representation.

Following extraction, the content undergoes structural analysis to identify the document's chapter and section hierarchy. This is represented as a tree structure in which each node corresponds to a heading level, forming the basis of the curriculum tree later displayed to the teacher. After the hierarchy is established, the content is divided into semantic chunks — cohesive text segments that correspond to a logical unit of knowledge, typically aligned with a subsection or section. Chunking is a critical design decision: chunks that are too large dilute retrieval specificity, while chunks that are too small lose necessary contextual information. A hybrid approach that respects the document's heading structure while enforcing a maximum token length is generally preferred.

In this project, teaching documents are processed through a pipeline that extracts raw content, detects heading structure, builds a curriculum tree, and segments the content into chunks with associated chapter identifiers. The resulting chunk corpus serves as the foundation for embedding generation and semantic retrieval in the subsequent stages of the exam generation workflow.

% 1.5.2
\subsubsection{Vector Database}

After documents are processed and segmented into chunks, each text unit is converted into a dense vector representation — an embedding — using a specialized embedding model. Each embedding encodes the semantic characteristics of its corresponding text segment in a high-dimensional vector space, such that text segments with similar meaning are positioned close to each other. This representation enables semantic retrieval, in which queries are compared against stored content based on meaning rather than lexical overlap.

A vector database is the infrastructure component responsible for storing, indexing, and querying these embedding vectors efficiently. When a query is submitted, it is first embedded using the same model used to index the document chunks, and then compared against stored vectors using a similarity metric. Cosine similarity — which measures the angle between two vectors and is insensitive to vector magnitude — is commonly used for this purpose.

To enable fast approximate nearest neighbor (ANN) search at scale, vector databases employ indexing algorithms such as HNSW (Hierarchical Navigable Small World) or IVF (Inverted File Index). HNSW, in particular, builds a multi-layer graph structure that allows the database to navigate to the most promising regions of the vector space at progressively finer granularity, achieving query speeds in the low milliseconds even across millions of vectors with negligible accuracy loss. These indexing strategies are what make semantic retrieval practical in production systems.

A notable design pattern in multi-document retrieval is the use of namespace-based isolation, in which vectors belonging to different chapters or documents are stored under separate namespaces within the same database index. This allows the system to restrict retrieval to the specific scope selected by the teacher without requiring separate database instances.

In this project, knowledge chunks are embedded using a dense embedding model (BAAI/bge-m3, 768 dimensions) and stored in Pinecone, with each document assigned a dedicated namespace. When the teacher selects chapters for exam generation, retrieval is confined to the corresponding namespace, ensuring that the generated content remains strictly within the intended scope.

% ------------------------------------------------------------
\subsection{Prompt Engineering}

Prompt engineering is the discipline of designing and optimizing the input prompts given to LLMs in order to elicit more accurate, consistent, and task-appropriate responses. Since LLMs generate outputs probabilistically based on their input context, the way a prompt is structured — including its instructions, format specifications, examples, and contextual information — has a direct and substantial effect on output quality. Effective prompt engineering is particularly critical in systems where the LLM must adhere to domain-specific constraints, produce structured output, or operate within a multi-stage pipeline.

Several prompt design strategies are commonly employed to improve LLM performance. Zero-shot prompting provides a direct instruction without examples, relying entirely on the model's pre-trained knowledge. Few-shot prompting augments the prompt with a small number of representative input-output examples, enabling the model to infer the desired pattern from context. Chain-of-thought prompting encourages the model to generate intermediate reasoning steps before producing its final answer, which has been shown to significantly improve performance on complex reasoning tasks.

In the context of educational content generation, prompt engineering serves as the mechanism through which domain knowledge, pedagogical constraints, and quality standards are communicated to the LLM. For example, exam blueprint prompts specify the desired Bloom's taxonomy distribution across chapters, ensuring that the generated outline reflects a balanced cognitive progression. Question generation prompts include instructions to produce distinct distractors, justify correct answers in a rubric, and assign an appropriate difficulty level. Validation prompts may invoke chain-of-thought reasoning to assess whether a generated question accurately reflects its assigned Bloom level and whether its distractors are plausible yet clearly incorrect.

In this project, prompt engineering is embedded throughout the agent prompts used within the LangGraph pipeline. Each agent is equipped with a system prompt that encodes its specific role, constraints, and output format requirements. Additionally, the system uses few-shot examples — drawn from the retrieved document context — to guide the model toward content that is grounded in the source material rather than inferred from general knowledge.

% ------------------------------------------------------------
\subsection{Web Development Technologies}

% 1.6.1
\subsubsection{Frontend Technologies}

The frontend of the system is built using Next.js, a React-based framework that provides server-side rendering (SSR), static site generation (SSG), and an integrated API layer through route handlers. Server-side rendering is particularly advantageous for an educational platform because it improves initial page load performance and supports SEO for public-facing content, while client-side interactivity is preserved for dynamic features such as real-time generation monitoring and HITL review forms. The App Router introduced in Next.js 16 enables the use of React Server Components, reducing the amount of JavaScript sent to the client and improving responsiveness.

React is used as the component model due to its declarative, component-based architecture, which facilitates the management of complex UI states. In the exam generation interface, components must handle several interdependent state dimensions simultaneously: the current pipeline stage, the evolving list of generated questions, WebSocket messages for live progress, user edits to individual questions, and the approval or rejection state at each HITL checkpoint. React's state management ecosystem — combined with optimistic UI updates — provides the necessary tools to handle this complexity without sacrificing code maintainability.

Tailwind CSS is employed as the utility-first styling framework. By co-locating styling with component logic and eliminating the overhead of maintaining a separate stylesheet, Tailwind accelerates UI development and ensures visual consistency across the application. Its design system is extended with a custom theme that supports both light and dark modes, accommodating different user preferences.

Real-time communication between the client and server is achieved through WebSocket connections. In the exam generation workflow, the server streams pipeline events — such as retrieval progress, blueprint readiness, and question generation updates — to the client as they occur, enabling a live generation viewer without requiring the client to poll repeatedly.

% 1.6.2
\subsubsection{Backend Technologies}

The backend is implemented in Python using FastAPI, an asynchronous web framework chosen for its native support for concurrent request handling and its tight integration with Python's AI and data science ecosystem. FastAPI's async/await model is particularly well-suited for handling multiple concurrent WebSocket connections — one per active generation session — without blocking server resources. Additionally, FastAPI provides automatic request validation through Pydantic models, interactive OpenAPI documentation, and native dependency injection, which together accelerate development and reduce the likelihood of runtime errors.

Python is selected as the primary language because it provides the most mature and well-supported libraries for AI workloads, including LangChain and LangGraph for agent orchestration, sentence-transformers for embedding generation, and WeasyPrint for document export. This ecosystem integration reduces the friction of bridging traditional web services with AI processing pipelines.

PostgreSQL serves as the primary relational database, storing structured data including user accounts, document metadata, exam configurations, and question banks. Its support for JSONB columns enables the flexible storage of semi-structured data — such as exam configurations, curriculum trees, and question JSON — without sacrificing the ability to query specific fields.

Redis is used for multiple supporting roles: rate limiting to protect the generation endpoint from abuse, session management for WebSocket connection state, and caching of frequently accessed document metadata. Its pub/sub mechanism also enables the real-time broadcasting of pipeline events to connected WebSocket clients.

Celery, a distributed task queue for Python, is used to offload long-running processing tasks — such as document parsing, embedding generation, and RAG retrieval — from the main FastAPI request path into background workers. This ensures that the API remains responsive during asynchronous operations and allows the system to scale worker instances independently of the web server.

Docker is used to containerize the application components, ensuring consistent behavior across development, staging, and production environments. Each major component — the FastAPI web server, Celery workers, PostgreSQL, Redis, and MinIO (S3-compatible object storage) — is defined as a separate service in a Docker Compose configuration, enabling straightforward local development and reproducible deployments.

\end{document}
